\documentclass[10pt,a4paper]{amsart}
\usepackage[margin=19mm]{geometry}
\usepackage{amsmath,amssymb,mathtools}
\usepackage[scr=rsfs,cal=euler]{mathalfa}
\usepackage{xfrac}
\usepackage[unicode,hidelinks,bookmarks=false]{hyperref}
\newcommand{\ie}{i.e.\ }
\newcommand{\cf}{cf.\ }
\newcommand{\resp}{respectively\ }
\newcommand{\Subsection}{\subsection}
\providecommand{\nobreakdash}{\nobreak\hbox{-}}
\setlength{\emergencystretch}{2em}
\multlinegap0pt
\usepackage{bm}
\usepackage{bbm}
\usepackage{dsfont}
\usepackage{xspace}
\usepackage{cancel}
\usepackage{mathtools}
\usepackage[shortlabels]{enumitem}
\usepackage{amsthm}
\makeatletter
\@ifundefined{theorem}{\newtheorem{theorem}{Theorem}}{}
\@ifundefined{lemma}{\newtheorem{lemma}[theorem]{Lemma}}{}
\makeatother
\title[Cellular waists of hyperbolic spaces]{Cellular waists of hyperbolic spaces}
\author{Grigori Avramidi, Thomas Delzant}
\date{Source: arXiv:2606.13585v1 (CC BY 4.0)}
\begin{document}
\maketitle

\noindent\emph{Source and licence.} Cellular waists of hyperbolic spaces. Version arXiv:2606.13585v1 (CC BY 4.0), distributed under the Creative Commons Attribution 4.0 International licence, as recorded in the arXiv licence metadata for this exact version and shown on the abstract page https://arxiv.org/abs/2606.13585v1. Full rights notice: licenses/arxiv-2606.13585-CC-BY-4.0.txt.\par\smallskip

\noindent\textit{Editorial scope.} Counted unit: the Lemma whose source label is \texttt{vanishlemma} in the article (source0.tex lines 701--718, source label \texttt{vanishlemma}), delivered together with the article's own complete original proof of it (source0.tex lines 744--818, one \texttt{proof} environment of 75 lines). No mathematics is written, repaired or re-derived: statement and proof are the article's own text line for line. The proof is detached from the statement in the source: the article states the result at lines 701--718 and prints its proof at lines 744--818, and the proof environment's own header names the statement, so the association is the article's own. Required context retained uncounted: leibnitzformula (lines 388--392); cyclelemma (lines 417--420). Every definition, display and label of those lines is retained unchanged. Omitted from this package and not redistributed: 1--387, 393--416, 421--700, 719--743, 819--990. Nothing inside a retained excerpt is omitted or altered: every retained body is delivered exactly as the article's lines are written, so no omission receipt of any kind accompanies this package. Citations: the retained excerpts carry no citation command at all, so no bibliography entry of the article is transcribed and no reference is redistributed. Reference adaptation: none was needed and none was made; every reference inside the delivered unit resolves inside it, so no unresolved reference or citation is produced. Printed slips kept literal: no printed word, symbol, hypothesis or number of the counted statement or of its proof was altered. Layout and class: the article's own class is \texttt{amsart} with packages including babel, amsmath, amssymb, comment, enumitem, amsthm, xcolor, amscd, graphicx, blindtext; the wrapper is the lane's pinned \texttt{amsart} editorial wrapper at 10pt on A4, which restates the theorem-like environments the retained text needs and adds the article's own macro definitions that the retained text needs (none), with extra wrapper packages (bm, bbm, dsfont, xspace, cancel, mathtools, enumitem). The delivered numbering is the unit's own. Assets: no retained span contains a figure environment, a graphics-inclusion command, a PDF-page inclusion, a table or a TikZ picture; no image file is copied into this package, the package declares no asset dependency and no image file is redistributed with it. Licence: version arXiv:2606.13585v1 of this article is distributed under the Creative Commons Attribution 4.0 International licence, as recorded in the arXiv licence metadata for this exact version and shown on the abstract page https://arxiv.org/abs/2606.13585v1. A full rights notice is delivered at licenses/arxiv-2606.13585-CC-BY-4.0.txt. Characters: every retained excerpt is pure ASCII, so the delivered engine, XeLaTeX with its default Latin Modern text font, reports no missing character.
\begin{equation}
\tag{*}
\label{leibnitzformula}
\partial(\sigma\times c)=\varphi(\partial\sigma\times c)+(-1)^{|\sigma|}\sigma\times\partial c
\end{equation}

\begin{lemma}
\label{cyclelemma}
If $|\tau|=|\sigma|-1$ and $c$ is a cycle, then $c_{\tau}$ is a cycle. 
\end{lemma}

\begin{lemma}
\label{vanishlemma}
The map 
$
F\partial(\sigma\times-)%:C_{d-|\sigma|}(X_{\sigma})\rightarrow C_{d-1}(BG).
$
satisfies the following properties.
\begin{enumerate}[label=$\mathrm{(\roman*)}$]
\item\label{cycle} 
If $d-|\sigma|=k$ and $c$ is a cycle in $C_{d-|\sigma|}(X_{\sigma})$, then $F\partial(\sigma\times c)=0$, 
\item\label{chainvanish} 
if $d-|\sigma|>k$ and $c$ is a chain in $C_{d-|\sigma|}(X_{\sigma})$, then $F\partial(\sigma\times c)=0$, 
\item\label{chain1} 
$\partial F\partial(\sigma\times-)=0$ for $d>1$, and
\item\label{chain2}
$\varepsilon F\partial(\sigma\times-)=0$ for $d=1$. 
\end{enumerate}
\end{lemma}

\begin{proof}[Proof of Lemma \ref{vanishlemma}]



%Choose such a lift and denote it by $F(\sigma\times-)$. Then 
%$$
%\partial F(\sigma\times-)=F(\partial(\sigma\times-)).
%$$
Recall the formula (\ref{leibnitzformula}): $$
\partial(\sigma\times c)=\varphi(\partial\sigma\times c)+(-1)^{|\sigma|}\sigma\times \partial c
$$ 
and the definition of the chains $c_{\tau}$ via $$\varphi(\partial\sigma\times c)=\sum_{|\tau|<|\sigma|}\tau\times c_{\tau}.
$$


Applying $F$ to the formula (\ref{leibnitzformula}) we get %for the differential of the product $\sigma\times c$, we get
$$
F\partial(\sigma\times c)=\sum_{|\tau|<|\sigma|-1}F(\tau\times c_{\tau})+\sum_{|\tau|=|\sigma|-1}F(\tau\times c_{\tau})+(-1)^{|\sigma|}F(\sigma\times\partial c). 
$$
We also know (by Lemma \ref{cyclelemma}) that the chains $c_{\tau}$ are cycles if $|\tau|=|\sigma|-1$. 

From this and the inductive hypothesis, we conclude that \begin{itemize}
\item 
$F\partial(\sigma\times c)$ vanishes if
$|c|=k$ and $c$ is a cycle: 
\begin{itemize}
\item 
$F(\tau\times c_{\tau})=0$ since $|\tau\times c_{\tau}|=d-1$ and 
\begin{itemize}
\item 
if $|\tau|=|\sigma|-1$ then $|c_{\tau}|=k$ and $c_{\tau}$ is a cycle by Lemma \ref{cyclelemma}, 
\item
if $|\tau|<|\sigma|-1$ then $|c_{\tau}|>k$. 
\end{itemize}
\item 
$F(\sigma\times\partial c)=0$ since $\partial c=0$.
\end{itemize}
\item 
$F\partial(\sigma\times c)$ vanishes if
$|c|>k$:
\begin{itemize}
\item 
$F(\tau\times c_{\tau})=0$ since $|\tau\times c_{\tau}|=d-1$ and $|c_{\tau}|>k$, and 
\item 
$F(\sigma\times\partial c)=0$ since $|\sigma\times\partial c|=d-1$ and either 
\begin{itemize}
\item 
$|\partial c|=k$ and $\partial c$ is a cycle, or
\item 
$|\partial c|>k$.
\end{itemize}
\end{itemize}
\end{itemize}
This establishes points \ref{cycle} and \ref{chainvanish} of the lemma. 

Finally, note that by the inductive hypothesis we have defined $F$ on $C_{\leq d-1}(X')$ and verified that it is a chain map on this complex. So, since $|\partial(\sigma\times c)|=d-1$, we have 
$$
\partial F\partial(\sigma\times c)=F\partial^2(\sigma\times c)=0
$$
for $d>1$ and 
$$
\varepsilon F\partial(\sigma\times c)=F\varepsilon\partial(\sigma\times c)=0
$$
for $d=1$. This establishes points \ref{chain1} and \ref{chain2} of the lemma.% \ref{vanishlemma}.

%Finally, for $d=1$ we must have $|\sigma|=1$ and $|c|=0$ and it is enough to check for $\sigma=[v,w]$ and $c$ a vertex (with a positive sign in $C_0(X_{[v,w]})$). Then $\partial([v,w]\times c)=v\times c_v-w\times c_w$. Hence $F(\partial(\sigma\times c))=\phi(c_v)-\phi(c_w)$. The augmentation map $\varepsilon$ sends this to zero since $c_v$ and $c_w$ are points in the $G$-cover of $X$ that can be connected by a path.
%\begin{remark}
%There is probably a simpler thing to say for this augmentation argument.
%\end{remark}
%Hence, since $BG$ is aspherical and $C_{d-|\sigma|}(X_{\sigma})$ is a free module, $F(\partial(\sigma\times-))$ lifts to $C_d(BG)$. 

%Finally, let $F(\sigma\times-)$ be a lift of $F(\partial(\sigma\times-))$ that vanishes on boundaries if $d-|\sigma|=k$ and vanishes on all chains if $d-|\sigma|>k$. Such $F$ satisfies the properties in the Proposition. This finishes the proof of the inductive step. 

%So, if $|c|=k$ and $c$ is a boundary, or if $|c|>k$, we may define the lift by $F(\sigma\times c)=0$. This finishes the proof. 
\end{proof}

\end{document}
